% Dreaming Itself — ALIFE 2026 preprint draft
% Build: latexmk -pdf PAPER.tex   (or pdflatex twice)
% Content mirrors PAPER.md; every number is reproducible from results/REPORT.md.
\documentclass[11pt]{article}

\usepackage[margin=1.1in]{geometry}
\usepackage{newtxtext,newtxmath}
\usepackage{microtype}
\usepackage{booktabs}
\usepackage{amsmath}
\usepackage{graphicx}
\usepackage{caption}
\usepackage{enumitem}
\usepackage[round]{natbib}
\usepackage{xcolor}
\usepackage[colorlinks=true,linkcolor=black,citecolor=blue!50!black,urlcolor=blue!50!black]{hyperref}

\captionsetup{font=small,labelfont=bf}
\setlist{itemsep=2pt,topsep=4pt}

\newcommand{\Wq}{\ensuremath{W_q}}
\newcommand{\dtcoord}{\ensuremath{dt_{\mathrm{coord}}}}

\title{\bfseries Dreaming Itself: Discrete Language Beats Continuous Embeddings at
Matched Bandwidth in a Physically Embodied, Causally Necessary
Communication Task --- and Diffusion Bodies Make It Worse}
\author{Xining Li}
\date{August 2026}

\begin{document}

\maketitle

\begin{center}
\begin{minipage}{0.9\textwidth}
\small\itshape
Prepared for ALIFE 2026. This is a technical report / preprint draft, not
a peer-reviewed publication. All results, code, and raw logs are in the
companion repository
\href{https://github.com/xiningli/dreaming-together}{\texttt{dreaming-together}};
every number in this paper is reproducible from
\texttt{results/REPORT.md}, generated by
\texttt{dreaming\_together/evaluation/make\_report.py} from
\texttt{results/grid.csv} and \texttt{results/grid\_seeds.csv}.
\end{minipage}
\end{center}

\begin{abstract}
\noindent
We ask whether structuring a learned inter-agent communication channel as
discrete tokens, rather than a continuous embedding, changes team
performance when the two channels are held to \emph{exactly} the same
information bandwidth (40 bits/frame). We test this in a 2v2 MuJoCo
tank-combat task engineered so that the channel is not merely useful but
demonstrably causally necessary: a shield-bearing agent's own body
occludes the information it needs to time a firing window, and silencing
its teammate's channel costs it dozens of win-rate points. Three
conditions cross a coordinator representation (discrete 32-token
language vs.\ a 5-dimensional 8-bit-quantized embedding) with a
per-agent policy class (feedforward vs.\ diffusion), holding channel
bandwidth, network update counts, and evaluation protocol fixed across
conditions. The headline comparison --- language beats embedding within
the diffusion-bodied policy class --- replicates across three independent
training seeds: it holds in 3 of 3 seeds. A second, unplanned finding
replicates equally cleanly and runs the opposite direction from our
pre-registered prediction: teams with diffusion bodies and an embedding
coordinator are not merely non-superior to language, they are robustly
the \emph{worst} condition, reversing below even the simpler
feedforward-bodied embedding team in 3 of 3 seeds. Several other
pre-registered predictions (optimal messaging rate, bits-per-win,
robustness at high message rate, window-timing proxy) are evaluated but
only on a single training seed each, and we report them with that
caveat rather than the confidence warranted by the replicated headline
result. We also report a methodological finding of independent
interest: direct policy-gradient fine-tuning of the diffusion policies
failed to scale to this 2-vs-2 task across more than a dozen attempts,
and training only converged once we froze the diffusion generator and
learned a lightweight critic-based selector over its samples --- a
``sample-and-select'' architecture that is, to our knowledge, an
underused fallback for diffusion-under-RL failures at moderate policy
horizons.
\end{abstract}

\section{Introduction}

Learned multi-agent systems increasingly develop their own internal
communication channels --- a continuous vector passed between value heads,
policy heads, or memory modules. Two questions about this channel are
usually conflated: \emph{does communication help}, and \emph{does the
specific representational form of the message matter}. Most
emergent-communication work addresses the first question, typically in
low-dimensional referential games where a channel's mere existence is
already the interesting result, and where ``does it help'' is hard to
separate from ``is the game even solvable without it''
\citep{lazaridou2017,foerster2016,mordatch2018}. We isolate the second
question by construction: bandwidth is asserted equal --- 40 bits/frame
--- between a discrete 32-token vocabulary (5 bits/token $\times$
8-token message) and a continuous 5-dimensional bottleneck quantized to
8 bits per dimension, at process startup, in code
(\texttt{assert\_bandwidth\_parity()}), not merely in the paper's prose.
And we place the question inside a physically embodied task where
communication is not incidental to performance but load-bearing by
construction: a shield-bearing agent's own hull physically occludes the
geometric information it needs to judge when a firing lane is briefly
open, so its teammate's message is the only channel through which that
information can reach it at all.

This design is also a small transplant of a question usually asked in
cultural-evolution and iterated-learning research --- why do
compositional, discrete languages emerge and persist under repeated
transmission, while unstructured signaling schemes do not
\citep{kirby2001,skyrms2010} --- into a
reinforcement-learning-from-scratch setting, where the ``generations''
are gradient steps rather than a population of learners passing a
protocol to naive successors. We include one direct probe of this
framing: a learnability/transmission test (P8) that behavior-clones a
fresh, randomly initialized listener on a small number of logged
(state, message, action) triples and measures how much team performance
the clone recovers, predicting that a compositional discrete protocol
survives this bottleneck better than an entangled continuous one.

\paragraph{Contributions.}
\begin{enumerate}
\item A physically embodied 2v2 combat task in which channel necessity is
  measured, not assumed: silencing the coordinator costs the
  fully-trained language-coordinated team 31--47 percentage points of
  win rate depending on seed (Table~\ref{tab:seeds}), and a trained-deaf
  baseline confirms the ablation understates how much a team can
  compensate once it has adapted to silence (\S\ref{sec:nc}).
\item A matched-bandwidth, two-policy-class comparison (feedforward and
  diffusion bodies, $\times 2$ coordinator representations, plus a
  pure-language condition) that separates the coordinator's
  representational format from the per-agent action-generation prior.
\item A 3-seed replication of the headline ordering claim --- the first
  multi-seed evidence in this line of experiments --- which overturns the
  single-seed pilot's apparent tie between the two weakest conditions
  and replaces it with a robust, reversed ordering (\S\ref{sec:p1}).
\item A working ``sample-and-select'' recipe for training diffusion
  action policies under on-policy RL at a scale where direct fine-tuning
  reproducibly failed, plus a pre-registered failure ledger and
  executable-gate development discipline that we believe generalizes
  beyond this project (\S\ref{sec:pipeline}, \S\ref{sec:methodology}).
\end{enumerate}

\section{Related work}

Emergent multi-agent communication has been studied via discrete
signaling channels trained through differentiable relaxations or
REINFORCE \citep{foerster2016,lazaridou2017}, and via continuous or
compositional protocols whose structure is analyzed after training
rather than fixed by construction \citep{mordatch2018,chaabouni2020}.
These studies typically vary channel architecture without holding
bandwidth exactly constant, and study referential or simple cooperative
games rather than physically embodied, continuous-control tasks.
Iterated-learning models \citep{kirby2001,skyrms2010} motivate the
prediction that compositional, discrete codes survive transmission
bottlenecks better than entangled continuous ones, a claim we test
directly (P8) rather than only motivationally. On the policy side,
diffusion models for action generation
\citep{ho2020,janner2022,chi2023} are normally trained by imitation,
not on-policy RL; we report what happened when we tried to close that
loop with PPO \citep{schulman2017} at 2v2 scale. The physics substrate
is MuJoCo \citep{todorov2012}.

\section{Task: 2v2 tank combat with a measured causally necessary channel}

Two teams of two wheeled tanks fight in an $8\,\mathrm{m} \times
8\,\mathrm{m}$ arena (MuJoCo rigid body physics, 2\,ms physics step,
50\,ms policy step). Each team has a fixed division of labor: a
\textbf{SHIELD} agent carries an occlusion-capable plate that blocks
incoming pellets and takes zero damage by construction, and a
\textbf{SHOTGUN} agent fires 8-pellet volleys (6\,HP each, 100\,HP
total, $8^\circ$ cone, $<3\,\mathrm{m}$ range, 1.2\,s cooldown). Both
agents share one 5-dimensional action interface --- two track
velocities, two 2-DOF arm joint targets (pan/tilt), and a trigger ---
uniform across roles and across every experimental condition, so no
condition gets a different action space. A team loses when either of
its tanks reaches 0\,HP; episodes cap at 30\,s (draw).

The firing window is emergent geometry, not a scripted event: \Wq, the
fraction of the shotgun's firing cone unobstructed by its own team's
shield, must exceed a threshold before firing is tactically sound, and
in the tight column formation the task rewards, the shield's own hull
generically occludes part of the shotgun's forward camera view of the
enemy. The shield agent therefore cannot, from its own observations
alone, reliably know when the shotgun's shot is both ready and safe to
take --- that information has to arrive from the coordinator. We do not
merely assert this: \S\ref{sec:nc} measures it, both as an eval-time
ablation and against a baseline trained from scratch with no channel at
all.

\section{Experimental design}

\subsection{Conditions and the bandwidth constraint}
\label{sec:conditions}

\begin{table}[ht]
\centering
\small
\begin{tabular}{@{}lll@{}}
\toprule
Condition & Per-agent policy & Team coordinator \\
\midrule
\textbf{A} & Feedforward MLP ($\sim$1M params, tanh-squashed, & Continuous embedding: 5-D bottleneck, \\
           & 8-step action horizon output)                    & 8-bit quantized per dim \\
\textbf{B} & Diffusion policy (DDPM, $\epsilon$-prediction,   & Continuous embedding (same as A) \\
           & 8-step DDIM at deployment)                       & \\
\textbf{C} & Diffusion policy (same as B)                     & Discrete language: 32-token vocabulary, \\
           &                                                  & 8-token messages, 4-layer Transformer \\
           &                                                  & decoder, $d=128$ \\
\bottomrule
\end{tabular}
\caption{Experimental conditions.}
\label{tab:conditions}
\end{table}

A isolates the coordinator representation with a simple body; B vs.\ C
isolates it with a diffusion body; A vs.\ B isolates the diffusion body
prior at fixed coordinator representation. Both coordinator channels are
asserted equal at 5 bits/token $\times$ 8 tokens $=$ 5-D $\times$ 8-bit
$=$ \textbf{40 bits/frame}, checked in code at every training and
evaluation entry point, not merely claimed in prose --- the vocabulary
size (32, not the originally sketched 40) was chosen specifically so
$\log_2(32) \times 8$ lands on an integer matching the embedding side
exactly. The coordinator emits a message every \dtcoord\ policy steps;
all reported stacks were \emph{trained} at $\dtcoord = 250$\,ms and
evaluated across a sweep of $\dtcoord \in \{50, 100, 250, 500, 1000,
2000\}$\,ms without per-rate fine-tuning.

Observations are a 16-D proprioceptive vector (heading, velocity, arm
state, implement state, game clock/HP) shared identically across roles
and conditions, plus role-specific privileged extras (nearest/second
enemy and teammate bearing, distance, and alive flags --- 11 additional
dimensions) used for the measured A/B/C/NC comparison reported here. A
segmentation-vision pipeline (single-channel $64 \times 64$ rendered
observations, 6 semantic classes, a pretrained encoder with a
bearing-prediction auxiliary) was built and separately validated to
camera-only performance parity with the privileged-state policies
(\S\ref{sec:methodology}, and see the repository's
\texttt{HANDOFF.md}), but \textbf{the measured comparison in this paper
uses the privileged 27-D state, not the pixel-based encoder} --- a scope
reduction from the original design that we flag explicitly as a
limitation (\S\ref{sec:limitations}): these results characterize
communication given accurate self/other state estimation, not
communication that must also compensate for a noisy learned visual
estimate of the same state.

\subsection{Evaluation protocol --- a deliberate deviation from self-play}
\label{sec:evalprotocol}

The original design specified self-play evaluation. We evaluate instead
against a single frozen, hand-calibrated scripted opponent
(\texttt{EliteScriptedTeam}), identical for every condition, rate, and
seed. This is a considered substitution, not an afterthought: a fixed
opponent removes co-evolving-opponent curricula as a confound, so a win
rate difference between conditions can be attributed to the
communication channel rather than to which condition's training
happened to face an easier or harder moving target. The cost is that
our win rate says nothing about \emph{inter-condition} head-to-head
play --- we do not know whether team C would beat team B directly, only
that both beat the same fixed yardstick at different rates. Every
reported cell is 500 episodes with bootstrap 95\% confidence intervals
(10{,}000 resamples).

\subsection{Training pipeline and the diffusion-under-RL failure}
\label{sec:pipeline}

Training proceeds in stages, each behind an executable pass/fail gate
(\S\ref{sec:methodology}): individual combat skills first
(behavior-cloned from a closed-form inverse-kinematics expert, then
PPO-refined, per role, no coordinator); then pair coordination, where a
scripted coordinator seeds the protocol through the \emph{actual}
message-passing interface the learned coordinator will later drive (so
the listener never has to adapt to a distribution shift between ``the
interface it learned'' and ``the interface it will be driven
through''), followed by coordinator PPO against frozen listeners, then
a joint fine-tune phase.

For the diffusion conditions (B, C), direct on-policy fine-tuning of the
diffusion action policy --- treating PPO as if the diffusion sampler
exposed a tractable log-probability, or substituting elite
self-imitation on top-return episodes for a policy gradient --- failed
to scale from the single-agent aim task (where it worked) to the full
2v2 task across more than a dozen attempted recipes, with several
distinguishable failure modes: naive exponentially-weighted advantage
reweighting on all self-generated samples degrades an already-competent
policy ($1.00 \rightarrow 0.35$ win rate on the aim-task pilot); a
multi-seed grid of six direct-fine-tuning attempts at 2v2 scale reached
at best 0.40 win rate; and joint coordinator/listener fine-tuning under
elite self-imitation catastrophically drifted a competent listener to
0.00 win rate when the coordinator's output distribution shifted under
it. The resolution, consistent with a pre-registered fallback in the
design document, is \textbf{sample-and-select}: freeze a diffusion
policy distilled from a certified feedforward teacher, draw $K=8$
action-horizon candidates from it per decision via 8-step DDIM
sampling, and train only a small feedforward critic
($\mathit{obs}, \mathit{action} \rightarrow \mathrm{scalar}$) by
return-to-go regression to score and select among the frozen
candidates. The generative body prior under test (A-vs-B) is untouched;
all the learning happens in a selector an order of magnitude smaller
than the generator itself. This architecture is what ``diffusion
policy'' denotes for every reported B and C number in this paper ---
not a directly RL-fine-tuned diffusion net.

\subsection{Multi-seed replication}
\label{sec:multiseed}

The original design specifies at least 3 independent training seeds per
condition. Our first complete measured pass reported single-seed
results only (seed 1). For this paper we trained two additional
independent seeds (2, 3) for every condition using the identical
recipe, and evaluated each at the trained rate (250\,ms), live and with
the channel zeroed, 500 episodes per cell. One implementation subtlety
mattered here and is worth recording for anyone reproducing this style
of experiment: seeds 2 and 3 initially reused seed 1's frozen
token-to-embedding table during a scripted-teacher distillation step,
which would have silently taught seed-2/3 listeners to interpret a
\emph{different} seed's vocabulary than the one their own coordinator
would eventually speak. We caught and fixed this before training the
reported seeds; it is the kind of bug that would not crash anything,
only quietly degrade a replication's validity.

\section{Results}

\subsection{Communication is causally necessary, and eval-time ablation
understates how much}
\label{sec:nc}

Table~\ref{tab:nc} (reproduced from \texttt{results/REPORT.md}) shows
win rate at the trained rate, live and with the coordinator's message
zeroed at evaluation time, alongside a \textbf{trained-deaf baseline
(NC)}: the same feedforward listener architecture and training budget
as condition A, but with the channel zeroed from the very first
training update, so it never learns to expect a message in the first
place.

\begin{table}[ht]
\centering
\small
\begin{tabular}{@{}lccc@{}}
\toprule
condition & live & zeroed (eval-time ablation) & trained deaf (NC) \\
\midrule
A (FF + embedding)        & 0.814 [0.780, 0.848] & 0.594 [0.552, 0.636] & --- \\
B (diffusion + embedding) & 0.694 [0.654, 0.734] & 0.528 [0.484, 0.572] & --- \\
C (diffusion + language)  & 0.916 [0.890, 0.940] & 0.424 [0.382, 0.468] & --- \\
NC (trained deaf)         & --- & --- & 0.734 [0.696, 0.772] \\
\bottomrule
\end{tabular}
\caption{$W(\text{condition}, 250\,\text{ms})$, single seed (seed 1),
500 episodes/cell, 95\% bootstrap CI.}
\label{tab:nc}
\end{table}

The eval-time ablation is a lower bound on how much a team \emph{could}
compensate for silence, not a measure of how much a team \emph{does}
compensate once trained for it: NC beats every condition's zeroed
column, sometimes by a wide margin (C zeroed reads 0.424, 31 points
below NC). Separating channel value (live $-$ NC) from
protocol-dependence penalty (zeroed $-$ NC): condition C's channel is
worth $+18.2$ win-rate points over a trained-deaf team ($p < 0.0001$,
two-proportion bootstrap), condition A's is worth $+8.0$ points
($p = 0.0013$), and condition B's channel is statistically
indistinguishable from having no channel at all trained-for-silence
($-4.0$ points, $p = 0.92$, underpowered --- not a confirmed null, just
not a confirmed effect at $n = 500$). This is one finding we report
with only single-seed confidence: the NC baseline was trained once, not
replicated across seeds.

\subsection{P1, replicated: language beats embedding; and a robust,
unplanned reversal}
\label{sec:p1}

The pre-registered prediction P1 was $W(C) > W(B) > W(A)$ at the
trained rate. We report this at two evidence tiers, because they are
not the same claim: the single-seed point estimate, and the 3-seed
replication.

\textbf{$C > B$ holds in 3 of 3 independent training seeds}
(Table~\ref{tab:seeds}) --- this is the result we are confident
asserting as a general finding of this experiment, not an artifact of
one training run. \textbf{$B > A$ does not hold in any of the 3 seeds}
--- every seed shows the \emph{opposite} ordering, $A > B$, which is
itself a clean, replicated, but unplanned result: the diffusion-bodied
team with a continuous embedding coordinator is the weakest of the
three conditions, not merely tied with the simpler feedforward team as
the single-seed pilot suggested.

\begin{table}[ht]
\centering
\small
\begin{tabular}{@{}lccc@{}}
\toprule
seed & A: win / drop & B: win / drop & C: win / drop \\
\midrule
1 & 0.814 [0.780, 0.848] / $+22.0$ & 0.694 [0.654, 0.734] / $+16.6$ & 0.916 [0.890, 0.940] / $+49.2$ \\
2 & 0.952 [0.932, 0.970] / $+29.8$ & 0.764 [0.728, 0.802] / $+17.4$ & 0.880 [0.852, 0.908] / $+85.6$ \\
3 & 0.854 [0.824, 0.884] / $+9.8$  & 0.722 [0.684, 0.760] / $-15.6$ & 0.738 [0.700, 0.776] / $+69.2$ \\
\bottomrule
\end{tabular}
\caption{Per-seed win rate at 250\,ms (live) and causal drop in
percentage points (live $-$ zeroed), 500 episodes/cell.}
\label{tab:seeds}
\end{table}

We deliberately do not pool these 500-episode-per-seed samples into one
large bootstrap. With only 3 training runs, the seed \emph{is} the unit
of replication --- each row is one independent draw from ``what
training this condition tends to produce,'' and the right summary of 3
draws is ordering consistency across them, not a p-value manufactured
by treating 1{,}500 episodes as 1{,}500 independent samples of one
underlying distribution when they are actually 3 samples of 500
correlated ones. By that standard: $C > B$ in 3/3 seeds; $A > B$ in 3/3
seeds.

We also disclose, rather than filter out, each seed's outcome against
the pre-registered gate that certifies a stack before it enters the
comparison (win rate $\geq 0.75$ vs.\ the frozen calibrated opponent,
message diversity $\geq 0.5\times$ a scripted reference, causal drop
$\geq 15$ points): 5 of 9 seed-stacks passed, 4 failed. Seed 3 is the
visibly weaker run across every condition, and B seed 3 is the one cell
in this entire study with a \emph{negative} causal drop ($-15.6$: this
trained team's zeroed win rate exceeded its live win rate) --- a
training run in which the learned protocol was actively unhelpful. We
report this number rather than discard the run and retrain until a
passing seed appears, because a paper reporting only gate-passing seeds
would systematically understate how unreliable training an
embedding-coordinated diffusion team actually is, which is itself part
of the finding.

\subsection{Single-seed predictions (report with visible hedging)}
\label{sec:singleseed}

The remaining predictions were evaluated only on seed 1. We report them
because they are informative, but the reader should not extend to them
the confidence that Table~\ref{tab:seeds}'s replicated ordering earns.

\begin{itemize}
\item \textbf{P2} (optimal message rate for C falls in $[250,
  1000]$\,ms): \textbf{not supported} at the point-estimate level ---
  the observed argmax across the rate sweep is 100\,ms, outside the
  predicted band, though confidence intervals across adjacent rates
  overlap substantially (Table 1 of the full report). Single seed only.
\item \textbf{P3} (B's optimal rate is faster than C's): \textbf{tied},
  not confirmed --- both point estimates argmax at 100\,ms. Single seed
  only.
\item \textbf{P4} (bits-per-win is lower for C than B, each at its own
  best rate): the point estimate supports this --- 2{,}813 bits/win for
  C vs.\ 5{,}746 for B, roughly half --- but this is a ratio of two
  single-seed point estimates with no propagated uncertainty, so we
  present it as suggestive rather than established.
\item \textbf{P6} (C and B converge at the fastest, 50\,ms, rate):
  \textbf{not supported} --- the gap between C (0.916) and B (0.672)
  persists at 50\,ms in the seed-1 data ($p < 0.0001$ on the pooled
  500-episode comparison; this p-value is legitimate as a within-seed
  statistic, it just does not tell us whether a different training seed
  would reproduce the same gap size).
\item \textbf{P7} (window-timing latency is minimized near the
  shotgun's natural fire-cycle period): we built only a coarse proxy
  (fraction of coordinator periods in which the firing window happens
  to be open), which is roughly flat (0.89--0.95) across the whole rate
  sweep for condition C and does not by itself resolve the prediction.
  The full cue-to-open latency instrument this prediction calls for was
  not built.
\end{itemize}

\subsection{P8: protocol learnability under a transmission bottleneck}
\label{sec:p8}

We behavior-cloned fresh, randomly initialized listener networks on $K$
logged (observation, message, action) episode pairs from each certified
stack's own live coordinator, then measured team win rate with those
cloned listeners playing under the \emph{original} coordinator (150
evaluation episodes). The prediction --- that the discrete-language
protocol survives a small-sample transmission bottleneck better than
the continuous-embedding one, with the gap widening as the bottleneck
tightens --- is supported at the point-estimate level: at $K=10$ logged
episodes, the language-condition clone recovers 0.907 team win rate
versus 0.667 for the embedding condition (a 0.240 gap); at $K=50$ the
gap narrows to 0.193 (0.847 vs.\ 0.653). This result, like
\S\ref{sec:singleseed}, is single-seed ($N=150$ evaluation episodes,
one probe per condition); we report the direction and its consistency
with the iterated-learning motivation in \S1, without claiming the
statistical weight of the replicated Table~\ref{tab:seeds} result.

\subsection{P5 --- not measured}
\label{sec:p5}

The design specified a $\Gamma$ instrument measuring self/world
discrimination via denoising-error differences between contact and
non-contact frames, predicted to differ for the diffusion-bodied
conditions (B, C) but not the feedforward one (A). We did not build
this instrument. We report this as a dropped prediction, not a deferred
or negative one --- there is no evidence for or against it in this
study, only an absence.

\section{Discussion}

The clearest, most defensible finding in this study is not the one we
set out to test most directly (P1's full three-way ordering) but a
sub-claim of it: \textbf{holding bandwidth exactly fixed, a discrete
32-token protocol produces a more competent coordinated team than a
continuous 5-D embedding when both drive diffusion-bodied agents, and
this holds across every independent training run we attempted.} The
second, unplanned but equally replicated finding is arguably more
interesting for a venue interested in the \emph{mechanisms} of
representation and evolution: pairing a diffusion body with a
continuous coordinator is not a graceful middle ground between ``simple
body, simple channel'' (A) and ``complex body, structured channel'' (C)
--- it is the worst combination in the design, consistently
underperforming even the simplest condition. One candidate mechanism,
observed but not formally isolated during bring-up (documented in
\texttt{results/STAGE2\_SUMMARY.md}): discrete tokens force a
listener's interpretation to snap to one of 32 fixed vocabulary points
regardless of small perturbations in the coordinator's output, while a
continuous bottleneck lets an imperfectly-trained coordinator drift
through nearby-but-wrong regions of a 5-dimensional manifold that a
frozen or slowly-adapting listener has no comparable discretization to
fall back on. If this mechanism is right, it predicts that
\emph{robustness to an imperfect speaker}, not raw expressive capacity,
is where discreteness earns its keep --- a reframing worth testing
directly in future work rather than inferring post hoc.

The G6 gate-failure pattern in \S\ref{sec:p1} is itself worth treating
as data, not noise to be trained away: 4 of 9 independently trained
seed-stacks failed our own pre-registered acceptance bar, concentrated
in exactly the condition (B) and exactly the seed (3) that the
replicated ordering flags as weakest. A single-seed report --- the
shape most experiments in this space are published as --- would have
had a roughly even chance of either passing or failing to notice this
instability, depending on which one seed happened to be run.

\section{Development methodology: gates, a failure ledger, and why we
think this generalizes}
\label{sec:methodology}

This project's predecessor (a humanoid variant, abandoned) trained for
several hours before anyone noticed the agents were falling at step
$\sim$10 and the reward gradient for an entire shaping term never
reached the policy. We responded by converting every prior failure into
a binding rule enforced by an executable gate before the next stage of
compute could begin: e.g., ``combat is possible'' is a unit test
(scripted oracle vs.\ scripted oracle must produce hits and non-draw
outcomes) that runs before any learning code does; every camera has a
rendered point-of-view unit test, because a prior version's shield
agent was discovered --- by watching a video after the fact --- to be
blind, occluded by its own equipment; every reward term has both a
value-correctness test and a gradient-flow test asserting a synthetic
nonzero advantage on that term alone changes policy, not just critic,
parameters. We list this not as a methods footnote but as a claim: for
embodied multi-agent RL research specifically, where a silently-wrong
observation or a zero-gradient reward term can burn a multi-hour run
without producing any error message at all, a pre-registered failure
ledger with executable gates is cheap insurance against the single most
common way these projects fail --- not bad algorithms, but bad plumbing
discovered too late to matter.

\section{Limitations}
\label{sec:limitations}

\begin{itemize}
\item \textbf{Single frozen opponent, not self-play.} Every win rate in
  this paper is against one calibrated scripted team, identical across
  every condition, rate, and seed. This is a deliberate choice
  (\S\ref{sec:evalprotocol}) that removes co-evolution as a confound at
  the cost of saying nothing about direct inter-condition play.
\item \textbf{Uneven replication depth.} Only the headline 250\,ms
  comparison (P1) has 3-seed replication. The rate sweep (P2, P3, P6,
  P7) and bits-per-win (P4) are seed-1 only; the learnability probe
  (P8) and the no-communication baseline (\S\ref{sec:nc}) are each a
  single training run. Extending full replication to the rate sweep
  would require retraining and re-evaluating the entire 6-rate grid for
  two additional seeds per condition --- a substantially larger compute
  budget than the fixed-rate replication reported here, and the next
  highest-value extension of this work.
\item \textbf{Privileged, not visual, observations.} The measured
  comparison uses ground-truth bearing/distance state, not the
  pixel-based perception pipeline that was separately built and
  validated to parity on individual-agent skills. These results speak
  to coordination given accurate state estimation, not to coordination
  that must also compensate for visual noise or occlusion of the
  \emph{observation} itself (as opposed to occlusion of the
  \emph{firing window}, which is the mechanic under test).
\item \textbf{P5 was not attempted}, and is not implicitly supported or
  refuted by anything in this paper.
\item \textbf{Diffusion conditions use sample-and-select}, not directly
  RL-fine-tuned diffusion sampling (\S\ref{sec:pipeline}). Results
  characterize this specific frozen-generator/learned-selector
  architecture; whether a directly fine-tunable diffusion policy at
  this scale would reproduce the same ordering is an open question this
  study cannot answer.
\item \textbf{Fixed vocabulary and coordinator period.} The 32-token
  vocabulary and the 250\,ms training rate were fixed design choices,
  not swept as independent variables in their own right.
\end{itemize}

\section{Conclusion}

At exactly matched information bandwidth, in a task engineered so that
communication is measurably load-bearing rather than merely available,
a discrete 32-token protocol produces better-coordinated teams than a
continuous 5-D embedding when both drive diffusion-bodied agents --- a
result that now rests on 3-seed replication, not a single training run.
An unplanned but equally robust second finding --- that diffusion
bodies paired with a continuous coordinator are the worst combination
tested, not a middle ground --- is arguably the more novel
contribution, and we have offered a candidate mechanism (discreteness
as a robustness property against an imperfectly-trained speaker) that
future work should test directly rather than infer from this result
alone. We have tried throughout to keep the paper's confidence
calibrated to the evidence actually collected: the headline ordering is
replicated and should be read as such; the rate-sweep and learnability
results are single-seed and should be read with the corresponding
caution; and one prediction (P5) was never measured and appears here
only as an honest absence.

\section*{Reproducibility}

Everything reported here is generated by code in the companion
repository. \texttt{python -m pytest tests/} (107 tests) exercises the
physics and environment layer;
\texttt{dreaming\_together/evaluation/stage3\_grid.py} and
\texttt{stage3\_seeds.py} produce \texttt{results/grid.csv} and
\texttt{results/grid\_seeds.csv} respectively;
\texttt{dreaming\_together/evaluation/make\_report.py} regenerates
\texttt{results/REPORT.md}, including every number in
Tables~\ref{tab:nc}--\ref{tab:seeds} and
\S\ref{sec:singleseed}--\ref{sec:p8}, from those two CSVs. Training
entry points, gate scripts, and the full failure ledger
(\S\ref{sec:methodology}, expanded) are in
\texttt{DESIGN\_OF\_EXPERIMENT.md}.

\begin{thebibliography}{12}

\bibitem[Chaabouni et al.(2020)]{chaabouni2020}
Chaabouni, R., Kharitonov, E., Bouchacourt, D., Dupoux, E., and Baroni,
M. (2020). Compositionality and generalization in emergent languages.
In \emph{Proceedings of ACL}.

\bibitem[Chi et al.(2023)]{chi2023}
Chi, C., Feng, S., Du, Y., Xu, Z., Cousineau, E., Burchfiel, B., and
Song, S. (2023). Diffusion policy: Visuomotor policy learning via
action diffusion. In \emph{Robotics: Science and Systems (RSS)}.

\bibitem[Foerster et al.(2016)]{foerster2016}
Foerster, J., Assael, Y.~M., de Freitas, N., and Whiteson, S. (2016).
Learning to communicate with deep multi-agent reinforcement learning.
In \emph{Advances in Neural Information Processing Systems (NeurIPS)}.

\bibitem[Ho et al.(2020)]{ho2020}
Ho, J., Jain, A., and Abbeel, P. (2020). Denoising diffusion
probabilistic models. In \emph{Advances in Neural Information
Processing Systems (NeurIPS)}.

\bibitem[Janner et al.(2022)]{janner2022}
Janner, M., Du, Y., Tenenbaum, J.~B., and Levine, S. (2022). Planning
with diffusion for flexible behavior synthesis. In \emph{Proceedings of
ICML}.

\bibitem[Kirby(2001)]{kirby2001}
Kirby, S. (2001). Spontaneous evolution of linguistic structure: an
iterated learning model of the emergence of regularity and
irregularity. \emph{IEEE Transactions on Evolutionary Computation},
5(2).

\bibitem[Lazaridou et al.(2017)]{lazaridou2017}
Lazaridou, A., Peysakhovich, A., and Baroni, M. (2017). Multi-agent
cooperation and the emergence of (natural) language. In
\emph{Proceedings of ICLR}.

\bibitem[Mordatch and Abbeel(2018)]{mordatch2018}
Mordatch, I. and Abbeel, P. (2018). Emergence of grounded compositional
language in multi-agent populations. In \emph{Proceedings of AAAI}.

\bibitem[Schulman et al.(2017)]{schulman2017}
Schulman, J., Wolski, F., Dhariwal, P., Radford, A., and Klimov, O.
(2017). Proximal policy optimization algorithms.
\emph{arXiv:1707.06347}.

\bibitem[Skyrms(2010)]{skyrms2010}
Skyrms, B. (2010). \emph{Signals: Evolution, Learning, and
Information}. Oxford University Press.

\bibitem[Todorov et al.(2012)]{todorov2012}
Todorov, E., Erez, T., and Tassa, Y. (2012). MuJoCo: A physics engine
for model-based control. In \emph{Proceedings of IROS}.

\end{thebibliography}

\section*{Acknowledgments}

Implementation, training infrastructure, and analysis code were
developed with AI-assisted pair programming (Claude Code, Anthropic).

\end{document}
